\documentclass[../paper.tex]{subfiles}

\begin{document}

In this section, we present the results of our experimental evaluation of the 18 algorithms with openly available implementations. This section is divided into three main parts. The first part presents results for the real-world dataset shown in Table~\ref{tab:real_world_instances}, and the instances from the 10th DIMACS implementation challenge shown in Table~\ref{tab:dimacs_instances}. The second part investigates the execution time and memory consumption, as well as the effect that parallel processing has on the algorithms that support it. And finally, the third part presents results for our syntactically generated graphs with varying amounts of noise applied to the graph structure and the node features.

\subsection{State-of-the-Art Comparisons}

Figure~\ref{fig:comparisons} shows the main results of our experimental comparison. The first result that stands out in terms of modularity optimization is that \textsc{VieClus} achieves the best result on \emph{all} instances when it finishes within the time and memory constraints. There are two graphs where \textsc{VieClus} exceeded the memory limit: uk-2002 and uk-2007-05, both from the DIMACS set. This result aligns with the results presented by the authors of \textsc{VieClus}, who excluded these two graphs in particular because they were too large for a memetic algorithm~\cite{vieclus2018}.

(TODO, paragraph on exact solutions).

For graphs with known ground-truth clusters, we observe a much more even split between the algorithms. However, the \textsc{MAGI} method stands out as the best-performing algorithm on this dataset, with consistently high F1 scores and the best worst-case performance. It is worth noting that most algorithms optimize modularity, while the ground-truth clusters could have relatively low modularity. For instance, the ground-truth clusters for the Citeseer graph have a modularity score of \numprint{0.54}, much lower than the highest recorded modularity by \textsc{VieClus} of \numprint{0.89}. In fact, all tested algorithms compute clusterings with higher modularity scores than the ground-truth clustering on this graph. One reason for this discrepancy is how the ground-truth clusters are defined. The Citeseer graph is a citation network in which vertices are scientific publications and edges represent citations. The ground-truth clusters represent categories of each publication. As Figure~\ref{fig:comparisons} would indicate, there is no guarantee that a particular publication favors citations within its own category, which would have aligned perfectly with the modularity objective.

There are significant differences between learning-based methods when applied to data without vertex features. Algorithms such as \textsc{UCoDe}, which was among the strongest learning-based algorithms on real-world instances, struggle to produce meaningful clusterings on DIMACS graphs. \textsc{DMoN} and \textsc{CommDGI} perform well despite being designed for graphs with node features and a known number of clusters. Unsurprisingly, \textsc{GNNS} performs best in this category, being the only learning-based algorithm designed for graphs without features.

(Streaming paragraph)

(Exact paragraph)

\subsection{Computational Efficiency}

Judging computational efficiency is made difficult by the wide variety of methods. Figure~\ref{fig:efficiency} shows the relation between clustering quality and computational time to highlight the main tendencies of the algorithms. Reading the figure from left to right, the streaming algorithms \textsc{Hollocou} and \textsc{CluStRE} are by far the fastest. Next are the two closely related algorithms \textsc{Leiden} and \textsc{Louvain}, both scoring near the top in both quality and execution time, only beaten by \textsc{VieClus} in terms of solution quality. Unsurprisingly, most learning-based methods take a long time to produce a clustering, but they also achieve high solution quality.

\subfile{../figures/time_quality.tex}

% \subfile{../tables/memory.tex}

\subsection{Syntetic Graphs}

To further evaluate the impact and sensitivity of graph features across clustering algorithms, we include experiments with syntactically generated graphs with vertex features. The generator we use is called the attributed, degree-corrected stochastic block model~\cite{adcsbm2022} (ADC-SBM). This generator plants $k$ clusters in the graph, and generates random edges via a distribution that is conditional on those clusters. It does so using the two input parameters $d$ and $d_{out} \leq d$, where $d$ is the expected degree for each vertex and $d_{out}$ is the expected number of neighbors for a node in a different cluster than its own. The ADC-SBM generator also produces node features that are correlated with the community structure. It does so using the two input parameters $\sigma$ for intra-cluster center variance and $\sigma_c$ for cluster center variance.

We conduct two sets of experiments with the ADC-SBM generator. First, we generate a sequence of graphs with gradually higher $d_{out}$, causing the graph structure to carry less information about the original $k$ clusters. And second, we generate graphs with strong community structure as we gradually increase the cluster center variance $\sigma_c$, leading the graph features to carry less information about the clusters.

% \subfile{../figures/adcsbm_plot.tex}

Figure~\ref{fig:adcsbm} shows the results of our two experiments with the ADC-SBM graph generator. Since this experiment aims to show how feature-aware methods respond to scenarios in which either the graph or the feature signal is distorted, we evaluate only methods that accept graph features as input. For reference, we also include \textsc{VieClus}, which performed best among the pure modularity-optimizing methods. When the feature signal is strong, \textsc{MAGI} and \textsc{DMoN} are the best performers, and at around $d_{out}=11$, all the feature-aware methods perform better than \textsc{VieClus}. It is also clear from this point of view that \textsc{MAGI} relies more on graph features than the graph structure, evident by the strong performance when the graph structure is distorted, and weak performance when the feature signal is distorted.

\textsc{DMoN} performs well in both scenarios, demonstrating its ability to capture community structure in both the graph and the features without relying too heavily on either one. Similarly, \textsc{CommDGI} displays similar performance, but is noticeably worse than \textsc{DMoN} across the board. \textsc{DGCluster} is somewhat opposite to \textsc{MAGI}, favoring the graph structure over the features. The same holds for \textsc{UCoDe}, but to a lesser extent.

\ifSubfilesClassLoaded{%
    \bibliography{../paper.bib}%
}{}

\end{document}